Als Kreisradius nehmen wir die Drahtlänge $r = \rHwO$ (der Abstand der Hände von der Drehachse und der Kugelradius werden vernachlässigt). Die Frequenz ist $f = \fHwF = \fHw$. In einer Umdrehung legt die Kugel den Umfang $2\pi r$ zurück:
\begin{align}
 v &= \vHwF = 2\pi\cdot\rHw\cdot\fHw = \vHw \approx \result{\vHwP{0}{2}}
\end{align}
Die Beschleunigung der gleichförmigen Kreisbewegung ist die Zentripetalbeschleunigung:
\begin{align}
 a &= \aHwF = \frac{\left(\vHw\right)^2}{\rHw} = \aHw \approx \result{\aHwP{0}{2}}
\end{align}
Das sind etwa sechzig Mal die Fallbeschleunigung. Die Masse der Kugel wird für diese beiden Grössen nicht gebraucht.
